\documentclass[a4paper]{article}
\usepackage[utf8]{inputenc}
\usepackage[russian,english]{babel}
\usepackage[T2A]{fontenc}
\usepackage[left=10mm, top=20mm, right=18mm, bottom=15mm, footskip=10mm]{geometry}
\usepackage{indentfirst}
\usepackage{amsmath,amssymb}
\usepackage[italicdiff]{physics}
\usepackage{graphicx}
\graphicspath{{images/}}
\DeclareGraphicsExtensions{.pdf,.png,.jpg}
\usepackage{wrapfig}
\usepackage{float}
\usepackage{subcaption}
\usepackage{enumitem}
\usepackage{caption}
\usepackage{array}
\captionsetup[figure]{name=Рисунок}
\captionsetup[table]{name=Таблица}

\title{\underline{Отчет о выполненой лабораторной работе 4.7.3}}
\author{Воронин Денис, Б04-407}

\begin{document}

\maketitle

\begin{center}
\textbf{\Large Изучение поляризованного света}
\end{center}

\textbf{Цель работы:} ознакомление с методами получения и анализа поляризованного света.

\textbf{В работе используются:} оптическая скамья с осветителем; зелёный светофильтр; два поляроида; чёрное зеркало; полированная эбонитовая пластинка; стопа стеклянных пластинок; слюдяные пластинки разной толщины; пластинки в 1/4 и 1/2 длины волны; пластинка
в одну длину волны для зелёного света (пластинка чувствительного оттенка).

\section{Аннотация}

По окончании работы должны будем: \par
с помощью чёрного зеркала определить разрешённые направления поляроидов; определить характер поляризации
света, прошедшего стопу и отражённого от неё под углом Брюстера;
оценить угол Брюстера для эбонита; выделить пластинки $\frac{\lambda }{2}$ и $\frac{\lambda }{4}$;
определить направления большей и меньшей скоростей для пластинки
$\frac{\lambda }{4}$; исследовать интерференцию поляризованных лучей.

\section{Теоретические сведения}

\subsubsection{Определение направления разрешённой плоскости колебаний поляроида}
	
Определить направление разрешённых колебаний поляроида проще всего с помощью чёрного зеркала. При падении света на отражающую поверхность под углом Брюстера свет в отражённом луче почти полностью поляризован, а вектор $\mathbf{E}$ параллелен отражающей поверхности. Луч света, прошедший поляроид и отразившийся от чёрного зеркала, имеет минимальную интенсивность при выполнении двух условий: во-первых, свет падает на отражающую поверхность под углом Брюстера и, во-вторых,
в падающем пучке вектор E лежит в плоскости падения.

Вращая поляроид вокруг направления луча и чёрное зеркало вокруг оси, перпендикулярной лучу, методом последовательных приближений можно добиться минимальной яркости луча, отражённого от зеркала, и таким образом определить разрешённое направление поляроида. Зная угол поворота зеркала (угол Брюстера), можно определить коэффициент преломления материала, из которого изготовлено
зеркало. Из определения угла Брюстера следует
\begin{equation} \label{BrewstersAngle}
    n = \tg{\varphi_\text{Б}}.
\end{equation}

\subsubsection{Получение эллиптически поляризованного света}

Эллиптически поляризованный свет можно получить из линейно поляризованного с
помощью двоякопреломляющих кристаллических пластинок.

Двоякопреломляющая пластинка имеет два взаимно перпендикулярных главных направления, совпадающих с осями эллипсоида диэлектрической проницаемости. Волны, поляризованные вдоль главных направлений, распространяются в пластинке с разными скоростями, не изменяя характера своей поляризации. Эти волны называются главными. Мы будем обозначать показатели преломления для главных волн через $ n_x $ и $ n_y $, где $ x $ и $ y $ --- главные направления кристаллической пластинки (рис. \ref{pic_1}).
\begin{figure}[!hb]
	\centering
	\includegraphics[width=0.3\linewidth]{p1.png}
	\caption{{ Разложение линейно поляризованного света по главным направлениям двоякопреломляющей пластинки}}
	\label{pic_1}
\end{figure}

Пусть на пластинку падает линейно поляризованная волна, электрический вектор которой ориентирован под некоторым углом $\alpha$ к оси
$x$. Разложим вектор $\mathbf{E}$ на составляющие $E_x$ и $E_y$. На входе пластинки $E_x$ и $E_y$ находятся в фазе. На выходе из-за разности скоростей между ними появляется разность хода $d(n_x - n_y)$, при этом сдвиг фаз определяется соотношением
\begin{equation}\label{PathDiff}
\Delta \varphi =  \dfrac{2\pi}{m} = k d(n_x - n_y)
\end{equation}

Как уже отмечалось, при сложении двух взаимно перпендикулярных колебаний, обладающих некоторым сдвигом фаз, образуется колебание, поляризованное по эллипсу.

\begin{figure}[!hb]
	\centering
	\includegraphics[width=0.35\linewidth]{p2.png}
	\caption{{Поворот направления колебаний с помощью пластинки в $\lambda/2$}}
	\label{pic_2}
\end{figure}

Рассмотрим практически важные частные случаи.

\begin{enumerate}
 		
\item Пластинка даёт сдвиг фаз $ 2\pi $ (пластинка в длину волны $ \lambda $). В результате сложения волн на выходе пластинки образуется линейно поляризованная волна с тем же направлением колебаний, что и в падающей волне.

\item Пластинка даёт сдвиг фаз $ \pi $ (пластинка в полдлины волны $ \lambda / 2 $). На выходе пластинки снова образуется линейно поляризованная волна. Направление $ bb' $ колебаний этой волны повёрнуто относительно направления $ aa' $ колебаний падающей волны (рис. 2). Как нетрудно сообразить, направление $ bb' $ является зеркальным отображением направления $ aa' $ относительно одного из главных направлений пластинки. Такую пластинку используют для поворота направления колебаний линейно поляризованного света.
	
\item Пластинка создаёт между колебаниями сдвиг фаз $ \pi/2 $ (пластинка в четверть длины волны). При сложении двух взаимно перпендикулярных колебаний, имеющих разность фаз $ \pi/2 $, образуется эллипс, главные оси которого cовпадают с координатными осями $ x $ и $ y $. При равенстве амплитуд возникает круговая поляризация.
 	
\end{enumerate}

Следует отметить, что, говоря о пластинках $ \lambda , \lambda/2, \lambda/4  $ и т. д., всегда подразумевают какую-либо вполне определённую монохроматическую
компоненту (например, пластинка $ \lambda/2 $ для зелёного света). Если на двоякопреломляющую пластинку падает не монохроматический свет, то на
выходе из неё для разных спектральных компонент эллипсы поляризации будут различными.

\subsubsection{Анализ эллиптически поляризованного света}

Анализ эллиптически поляризованного света сводится к нахождению главных осей
эллипса поляризации и к определению направления вращения электрического вектора.

Главные оси эллипса поляризации определяются с помощью анализатора по максимуму и минимуму интенсивности проходящего света. Направление вращения электрического вектора может быть найдено с помощью пластинки в четверть длины волны, для которой известно, какая из главных волн, $ E_x $ или $ E_y $, имеет б\'{o}льшую скорость распространения (и соответственно меньшее значение показателя преломления).

Выберем для определённости координатные оси x и y на пластинке так, чтобы $ nx < ny $. В этом случае главная волна $ E_x $ имеет большую скорость распространения. Поместим такую пластинку на пути эллиптически поляризованного света и совместим главные направления пластинки $ \lambda/4 $ с главными осями эллипса поляризации. На выходе из этой пластинки сдвиг фаз между $ E_x $ и $ E_y $ вместо $ \pi/2 $ станет равным нулю или $ \pi $. Свет окажется линейно поляризованным. Из двух возможных значений сдвига фаз, 0 или $ \pi $, реализуется одно: то, которое соответствует имеющемуся в волне направлению вращения электрического вектора.

Рассмотрим, например, случай, когда электрический вектор в эллиптически поляризованной волне вращается против часовой стрелки, если смотреть навстречу лучу. В этом случае, очевидно, в волне, падающей на пластинку в $ \lambda/4 $, колебание $ E_y $ отстаёт по фазе на $ \pi/2 $ от колебания $ E_x $. При прохождении через пластинку разность фаз увеличивается до $ \pi $. Таким образом на выходе из пластинки возникают линейно поляризованные волны со сдвигом фаз $ \pi $. Сложение этих волн даёт плоскополяризованную волну, электрический вектор которой располагается во втором и четвёртом квадрантах координатной системы $x, y$.

Рассуждая аналогичным образом, найдём, что при вращении электрического вектора по часовой стрелке направление колебаний в линейно поляризованной волне, выходящей из пластинки, располагается в первом и третьем квадрантах. Определяя направление колебаний на выходе из пластинки с помощью поляроида, можно, таким образом, определить характер эллиптической поляризации (вращение против или по часовой стрелке).

\subsubsection{Пластинка чувствительного оттенка}

\begin{figure}[!hb]
	\centering
	\includegraphics[width=0.2\linewidth]{p3.png}
	\caption{{Пластинка чувствительного оттенка}}
	\label{pic_3}
\end{figure}

Выше предполагалось известным, какому из двух главных направлений пластинки в четверть длины волны соответствует большая скорость распространения света.
Установить это можно различными способами, например с помощью
пластинки чувствительного оттенка (так называют пластинку в $\lambda$
для зелёной спектральной компоненты, $\lambda = 560$ нм).

Пластинка имеет форму стрелы (рис. 3), вдоль оси которой расположено главное направление, соответствующее большей скорости распространения.

Если между скрещенными поляроидами поместить пластинку чувствительного оттенка
($ \lambda $) и пластинку в $ \lambda/4 $ так, чтобы их главные
направления совпадали, цвет пластинки изменится. Если у пластинки чувствительного оттенка и пластинки в $ \lambda/4  $совпадут главные направления, соответствующие большей скорости распространения, то разность хода между $ E_x $ и $ E_y $ для зелёного света составит уже $ 5\lambda/4 $. Это соответствует разности хода в $ \lambda $ для света с большей длиной волны, т. е. для "<более красного"> света. При освещении
этих пластинок (напомним, что они расположены между скрещенными поляроидами) белым светом теперь погасится не зелёная, а красная
часть спектра, и проходящий свет будет казаться зеленовато-голубым.
Если же главные направления, соответствующие большей скорости распространения, у пластинки чувствительного оттенка и у пластинки
в $ \lambda/4 $ окажутся перпендикулярными, то проходящий свет приобретёт
оранжево-желтую окраску (погасится фиолетово-голубая часть спектра).

Изменение цвета позволяет, таким образом, определить, какое из
главных направлений пластинки в $ \lambda/4 $ соответствует большей скорости
распространения.

\subsubsection{Интерференция поляризованных лучей}

\begin{wrapfigure}{l}{0.3\textwidth}
  \centering
  \includegraphics[width=0.2\textwidth]{p4.png}
  \caption{К объяснению интерференции поляризованных лучей
}\label{fig:example}
\end{wrapfigure}

Тонкие двоякопреломляющие пластинки, помещённые между поляроидами, кажутся окрашенными. Эта окраска может быть истолкована как результат интерференции поляризованных лучей. На рис. 4 представлена схема для
случая скрещенных поляроидов.

Здесь $ p_1p'_1 $ --- разрешённое направление колебаний поляризатора
(первого поляроида); $ x, y $ --- координатная система, связанная с главны-
ми направлениями двоякопреломляющей пластинки; $ p_2p'_2 $ --- разрешённое направление колебаний анализатора (второго поляроида). Волны
$ E_x  $ и $ E_y $ на выходе из пластинки когерентны, но не могут интерферировать, так как $ E_x \perp  E_y $. Волны $ E_1 $ и $ E_2 $ на выходе второго поляроида
также являются когерентными и к тому же поляризованы в одной плоскости. Эти волны интерферируют между собой. Результат интерференции определяется зависящим от длины волны сдвигом фаз между $ E_1 $
и $ E_2 $. В результате интерференции поляризованных лучей пластинка, освещаемая белым светом, кажется окрашенной.

Если поворачивать двоякопреломляющую пластинку, расположенную между
скрещенными поляроидами, то соотношение амплитуд волн $ E_1 $ и $ E_2 $ и разность фаз между ними не изменяются. Это означает, что цвет пластинки при её поворотах не меняется, а меняется только интенсивность света. За один оборот пластинки интенсивность четыре раза обращается в нуль --- это происходит при совпадении главных направлений
$ x $ и $ y $ с разрешёнными направлениями колебаний поляроидов.

Если же двоякопреломляющую пластинку оставить неподвижной, а
второй поляроид повернуть так, чтобы разрешённые направления $ p_1p'_1 $
и $ p_2p'_2 $ совпали, то волны $ E_1 $ и $ E_2 $ приобретают дополнительный фазовый сдвиг на $ \pi $ для всех спектральных компонент; при этом их амплитуды изменятся так, что цвет пластинки изменится на дополнительный. 

\section{Ход работы}

\subsection{Определение разрешенных направлений поляроидов}

Соберем установку так, как написано в описании. Плоскость падения будет горизонтальна: 

\begin{figure}[!hb]
	\centering
	\includegraphics[width=0.4\linewidth]{p5.png}
	\caption{Определение разрешенных направлений поляроида}
	\label{pic_3}
\end{figure}

Поворачивая поляроид вокруг направления луча добьемся наименьшей яркости отраженного пятна. Это будет под углом Брюстера. Т.к. вектор E лежит в плоскости падения, то найдя таким образом минимум интенсивности мы определили разрешенное направление первого
поляроида (с маленьким пятном): $\textbf{горизонтальное}$. Отсчет по лимбу 293.
\par
Определим разрешенное направление второго поляроида с помощью первого. Для этого поставим поляроиды друг за другом и вращая второй поляроид добъемся исчезновение света.
Это будет означать, что мы скрестили их разрешенные направления и определили разрешенное направление второго поляроида: $\textbf{вертикальное}$. Отсчет по лимбу 215.

\subsection{Определение показателя преломления (угла Брюстера) для Эбонита}

Соберем установку: вместо черного зеркала поставим на скамью эбонитовую пластинку с круговой шкалой. Отметим по лимбу положение, соответствующее положению, которое перпендикулярно лучу: $\textbf{183}$.\par

Установив направление разрешенных колебаний поляроида P1 горизонатльно найдем угол поворота эбонита $\varphi_{\text{Б}}$ при котором интенсивность отраженного луча минимальна: $\textbf{123}$. Неточность установки составляет 0.5 градуса. После добавления светофильтра результат изменился незначительно. Это связано с тем, что так как 
Эбонит это диэлектрик, то по формулам френеля имеем разные амплитудные коэффициенты для поляризций. \par

Рассчитаем показатель преломления Эбонита: 
\[tg\varphi_{\text{Б}} = n \implies n = 1.55 \pm 5.02\]

Табличное значение составлет 1,66. 

\subsection{Исследование поляризации света в преломленном и отраженном от стопы лучах}

Поставим вместо Эбонита стопу и подберем для нее угол Брюстера. 

\newpage

\begin{figure}[!hb]
	\centering
	\includegraphics[width=0.4\linewidth]{p6.png}
	\caption{Исследование стопы}
	\label{pic_3}
\end{figure}

Осветим стопу неполяризованным светом, и рассматривая в поляроиды отраженнный и преломленный свет определим тип поляризации:
$\textbf{отраженный свет: вертикальная, преломленный свет: горизонтальная}$.

\subsection{Определение главных направлений двоякопреломляющих пластин}

\begin{figure}[!hb]
	\centering
	\includegraphics[width=0.4\linewidth]{p7.png}
	\caption{Определение главных направлений в пластинах}
	\label{pic_3}
\end{figure}

Поставим кристаллическую пластинку между скрещенными поляризаторами $P_1$ и $P_2$. Вращая пластинку вокруг направления луча и наблюдая за интенсивностью света, проходящего сквозь второй поляризатор, определим, при каком условии главные направления пластинки совпадают с разрешёнными направлениями поляризаторов. Минимумы и максимумы интенсивности чередуются через $45^\circ$, главные плоскости пластин совпадают с разрешёнными направлениями поляризаторов. Минимумы: для первой $\alpha_0 = 83$, для второй $\alpha_0 = 22$.

\subsection{Выделение пластин $\frac{\lambda }{2}$ и $\frac{\lambda }{4}$}

\begin{figure}[!hb]
	\centering
	\includegraphics[width=0.5\linewidth]{p1.jpg}
	\caption{Определение пластинок}
	\label{pic_3}
\end{figure}

\newpage

Для выделения пласт $\lambda/2$, $\lambda/4$ добавим к схеме, изображённой на рис.~2, зелёный фильтр и установим разрешённое направление первого поляризатора горизонтально, а главные направления исследуемой пластинки — под углом $45^\circ$ к горизонтали. Для первой получаем, что после поворота свет становится оранжевым, следовательно линейная поляризация и пластинка $\lambda/2$. Для второй: свет становится зелёным, следовательно круговая поляризация и пластинка $\lambda/4$.

\subsection{Определение быстрой и медленной оси в пластинке $ \lambda / 4 $}

\begin{figure}[!hb]
	\centering
	\includegraphics[width=0.4\linewidth]{p8.png}
	\caption{Определение направлений большей и меньшей скоростей}
	\label{pic_3}
\end{figure}

Определим «быструю» и «медленную» оси в пластинке $ \lambda / 4 $.
\\
Поставим между скрещенными поляроидами пластинку чувствительного оттенка, имеющую вид стрелки, и убедимся, что эта пластинка не меняет поляризацию зелёного света. Уберем зелёный фильтр и убедимся, что стрелка имеет пурпурный цвет. Это объясняется тем, что зелёная компонента линейно поляризованного света при прохождении пластинки не меняет поляризации и задерживается вторым поляроидом.
\\
Добавим к схеме пластинку $ \lambda / 4 $, главные направления которой совпадают с главными направлениями пластины $ \lambda $ и ориентированы под углом 45° к разрешённым направлениям скрещенных поляроидов. При повороте рейтера со стрелкой на 180° вокруг вертикальной оси цвет стрелки меняется от зелёно-голубого до оранжево-жёлтого. В первом случае у нас «быстрая» ось (они совпадают), во втором -- медленная.

\subsection{Интерференция поляризованных лучей}

Расположим между скрещенными поляроидами мозаичную слюдяную пластинку. Она собрана из 4-х узких полосок слюды, лежащих по сторонам квадрата (две полоски «толщиной» $ \lambda / 4 $ и но одной — $ \lambda / 2 $ и $ 3 \lambda / 4 $). В центральном квадратике слюды нет. Главные направления всех пластинок ориентированы параллельно сторонам квадрата.
\\\\
Вращая пластинку, пронаблюдаем за изменениями в отдельном квадратике. У нас изменяется интенсивность.
\\\\
Не трогая пластинки, повращаем второй поляроид. См. видео. 

\subsection{Определение направления вращения светового вектора в эллиптически поляризованной волне}

Снова поставим зеленый фильтр, а за ним между скрещенными поляроидами -- пластинку произвольной толщины $( \lambda / 4 )$.
\\\\
Получим эллиптически-поляризованный свет. Для этого установим разрешённое направление первого поляроида под углом 10°--20° к горизонтали так, чтобы вектор \textbf{E} падающего на пластинку света был расположен в первом квадранте. Установим разрешённое направление второго поляроида вертикально и, вращая пластинку, найдем минимальную интенсивность света, прошедшего второй поляроид. Вращая второй поляроид, убедимся, что свет поляризован эллиптически, а нс линейно. Таким образом, получим эллипс поляризации с вертикально ориентированной малой осью.
\\\\
Для определения направления вращения светового вектора в эллипсе установим между поляроидами дополнительную пластинку $ \lambda / 4 $ с известными направлениями «быстрой» и «медленной» осей, ориентированными по осям эллипса поляризации анализируемого света. В этом случае вектор \textbf{E} на выходе будет таким, как если бы свет прошел две пластинки $ \lambda / 4 $: свет на выходе из второй пластинки будет линейно поляризован. Если пластинки поодиночке дают эллипсы, вращающиеся в разные стороны, то поставленные друг за другом, они скомпенсируют разность фаз, и вектор \textbf{E} на выходе останется в первом и третьем квадрантах. Если же световой вектор перешёл в смежные квадранты, значит, эллипсы вращаются в одну сторону.
\\\\
После второго поляроида интенсивность света максимальна в первом квадранте. Значит, две пластины компенсируют друг друга, световой вектор остался в том же квадранте, эллипсы вращаются в разные сторону.

\section{Вывод}

Ознакомился с методами получения и анализа поляризованного света. С помощье черного зеркала определил разрешенные направления поляроиди. Определил характер поляризации света, прошедшего через стопу и отраженного от нее под углом Брюстера. Оценил угол Брюстера для Эбонита. Выделил пластинки $  \lambda / 2 $и$ 3 \lambda / 4 $. Определил направления большей и меньшей скоростей для пластинки $ \lambda / 4 $. Исследовал интерференцию поляризованных лучей.

\end{document}